\documentclass[11pt,a4paper]{article}

\usepackage[utf8]{inputenc}
\usepackage[T1]{fontenc}
\usepackage[portuguese]{babel}
\usepackage[a4paper,margin=2.2cm]{geometry}
\usepackage{hyperref}
\usepackage{xcolor}
\usepackage{enumitem}
\usepackage{longtable}
\usepackage{array}
\usepackage{titlesec}

\hypersetup{
  colorlinks=true,
  linkcolor=blue,
  urlcolor=blue,
  pdftitle={RunX Fase 2 - Guia de Apresentacao},
  pdfauthor={Grupo 16}
}

\setlist[itemize]{topsep=2pt,itemsep=2pt,parsep=0pt}
\setlist[enumerate]{topsep=2pt,itemsep=2pt,parsep=0pt}
\setlength{\parskip}{6pt}
\setlength{\parindent}{0pt}

\titleformat{\section}{\large\bfseries}{\thesection}{0.6em}{}
\titleformat{\subsection}{\normalsize\bfseries}{\thesubsection}{0.6em}{}

\begin{document}

{\LARGE \textbf{RunX Fase 2}}\\[4pt]
{\large Guia de apresentacao oral e funcional para demonstracao ao professor}\\[10pt]
\textbf{Projeto:} Backend Node.js + Express + Prisma + PostgreSQL\\
\textbf{Data:} 16/04/2026

\vspace{8pt}
\hrule
\vspace{8pt}

\section{Objetivo da demonstracao}

Mostrar que o backend esta \textbf{a funcionar fim a fim}, com:

\begin{itemize}
  \item autenticacao JWT;
  \item gestao de corridas, metas semanais e feed social;
  \item grupos e eventos (funcionalidade-chave da fase 2);
  \item validacao de dados e tratamento de erros;
  \item documentacao via Swagger.
\end{itemize}

\textbf{Mensagem principal para abrir a apresentacao:}

\textit{"Nesta fase entregamos um backend completo, com arquitetura por camadas, persistencia em PostgreSQL via Prisma, autenticacao segura e modulos sociais e de grupos operacionais."}

\section{Checklist antes da aula (10-15 min)}

\subsection{Checklist tecnico}

\begin{enumerate}
  \item Confirmar Node.js e npm instalados.
  \item Confirmar PostgreSQL ativo.
  \item Confirmar ficheiro \texttt{.env} presente e valido.
  \item Confirmar compilacao sem erros.
  \item Confirmar API a responder em \texttt{/health}.
\end{enumerate}

\subsection{Comandos de preparacao}

Executar na raiz do projeto:

\begin{verbatim}
npm install
npm run prisma:generate
npm run prisma:migrate
npm run build
PORT=3001 npm run dev
\end{verbatim}

\textbf{Nota importante:} usar porta \texttt{3001} evita conflitos caso o \texttt{3000} esteja ocupado por outro projeto local.

Teste rapido:

\begin{verbatim}
curl http://localhost:3001/health
\end{verbatim}

Esperado:

\begin{verbatim}
{"ok":true,"service":"runx-fase2-api"}
\end{verbatim}

\section{Roteiro de demo ao vivo (12 minutos)}

\begin{longtable}{|p{0.11\textwidth}|p{0.39\textwidth}|p{0.43\textwidth}|}
\hline
\textbf{Tempo} & \textbf{Acao funcional} & \textbf{Fala sugerida (guiao oral)} \\
\hline
0:00--1:00 & Mostrar contexto no README e stack do projeto. & "O backend foi organizado por camadas: rotas, controllers, services, models e schemas de validacao." \\
\hline
1:00--2:00 & Abrir \texttt{/health} e \texttt{/api-docs}. & "Aqui validamos rapidamente disponibilidade da API e documentacao interativa." \\
\hline
2:00--4:00 & Executar \texttt{/auth/register}, \texttt{/auth/login}, \texttt{/users/me}. & "A autenticacao usa JWT e os endpoints protegidos exigem Bearer token." \\
\hline
4:00--6:00 & Mostrar \texttt{/themes}, criar corrida em \texttt{/runs}, listar em \texttt{/runs}. & "As corridas ficam persistidas em PostgreSQL via Prisma, com validacoes de entrada." \\
\hline
6:00--7:00 & Mostrar \texttt{/weekly-goals/current} e \texttt{PUT /weekly-goals}. & "Acompanham-se metas por semana e o progresso e recalculado com base nas corridas." \\
\hline
7:00--8:30 & Mostrar \texttt{/posts}, \texttt{/feed}, reacoes e comentarios. & "Temos o modulo social completo: publicacoes, reacoes e comentarios." \\
\hline
8:30--10:30 & Criar grupo, criar evento, entrar no evento. & "Este e o destaque da fase 2: grupos, gestao de membros e eventos com participacao." \\
\hline
10:30--11:30 & Enviar pedido invalido para forcar erro 400. & "A API valida dados com Zod e devolve erros consistentes e explicitos." \\
\hline
11:30--12:00 & Fecho + referir relatorio PDF. & "Concluimos os objetivos da fase 2 e deixamos pontos de extensao para fases seguintes." \\
\hline
\end{longtable}

\section{Passos funcionais detalhados (Swagger ou curl)}

\subsection{1) Autenticacao}

\begin{verbatim}
POST /auth/register
POST /auth/login
GET  /users/me
\end{verbatim}

Guardar \texttt{accessToken} para os passos seguintes.

\subsection{2) Corridas e progresso semanal}

\begin{verbatim}
GET  /themes
POST /runs
GET  /runs
GET  /runs/week
\end{verbatim}

\subsection{3) Metas semanais}

\begin{verbatim}
PUT /weekly-goals
GET /weekly-goals/current
\end{verbatim}

\subsection{4) Feed social}

\begin{verbatim}
POST /posts
POST /posts/{id}/reactions
POST /posts/{id}/comments
GET  /feed
\end{verbatim}

\subsection{5) Grupos e eventos (fase 2)}

\begin{verbatim}
POST /groups
POST /groups/{id}/events
POST /groups/{id}/events/{eid}/join
GET  /groups/{id}/events
\end{verbatim}

\section{Plano B (se houver problemas durante a demo)}

\subsection{Porta ocupada}

Se a API nao arrancar no 3000:

\begin{verbatim}
PORT=3001 npm run dev
\end{verbatim}

\subsection{Erro de base de dados}

\begin{verbatim}
npm run prisma:generate
npm run prisma:migrate
\end{verbatim}

E confirmar \texttt{DATABASE\_URL} no \texttt{.env}.

\subsection{Token expirado}

Executar novo login e repetir pedidos com novo \texttt{accessToken}.

\subsection{Sem internet para APIs externas}

Nao bloqueia a demonstracao: integracoes externas desta fase estao opcionais/stub.

\section{Perguntas provaveis do professor (respostas curtas)}

\textbf{Porque nao existe homepage em / ?}\\
A API e orientada a recursos; as rotas publicas principais sao \texttt{/health} e \texttt{/api-docs}. Caminhos nao definidos devolvem 404 controlado.

\textbf{Como garantem qualidade dos dados?}\\
Validacao com Zod nos schemas, middlewares de validacao e codigos de erro consistentes.

\textbf{O que ficou para extensao futura?}\\
Sugestao de rota por IA e integracoes mais avancadas, mantendo endpoint e estrutura preparados.

\textbf{Como empacotaram a entrega?}\\
ZIP final sem \texttt{node\_modules} nem \texttt{.git}, incluindo codigo-fonte e ficheiros obrigatorios (incluindo \texttt{.env}).

\section{Fecho recomendado (30 segundos)}

\textit{"Concluimos os objetivos da fase 2 com backend funcional, persistente e documentado. A demonstracao mostrou autenticacao, corridas, metas, feed e grupos/eventos em funcionamento, com validacao robusta e estrutura pronta para evolucao."}

\vfill
\textbf{Links de demonstracao:}\
\begin{itemize}
  \item API: \url{http://localhost:3001}
  \item Health: \url{http://localhost:3001/health}
  \item Swagger: \url{http://localhost:3001/api-docs}
\end{itemize}

\end{document}
